\documentclass[11pt]{article}
\usepackage{setspace}
\usepackage{geometry}
\usepackage[hidelinks]{hyperref}
\geometry{margin=1in}
\setstretch{1.12}

\begin{document}

% --- Applicant Information ---
\noindent
\textbf{Zhongli Wang} \\
Bielefeld, Germany \\
Phone: +49 160 9541 6015 \\
Email: \href{mailto:nkwangzhongli@gmail.com}{nkwangzhongli@gmail.com} \\
Website: \href{https://sites.google.com/view/zhongli-wang}{https://sites.google.com/view/zhongli-wang}

\vspace{0.45cm}

% --- Recipient Information ---
\noindent
\textbf{Recruitment Committee} \\
Lingnan-60 Global Talent Recruitment (Post Ref.: 24/999) \\
Human Resources Office, Lingnan University \\
Tuen Mun, New Territories, Hong Kong SAR

\vspace{0.45cm}

% --- Subject Line ---
\noindent
\textbf{Application for Tenure-Track Assistant Professor (Computational Social Science) --- Post Ref.: 24/999}

\vspace{0.25cm}

\noindent Dear Members of the Recruitment Committee,\\[0.2cm]

I am writing to apply for the tenure-track Assistant Professor position under the Lingnan-60 Global Talent Recruitment (Post Ref.: 24/999). I am a PhD candidate in Economics at Bielefeld University and Universit\`a Cattolica del Sacro Cuore through the Marie Sk\l odowska--Curie EPOC programme (Grant No.\ 956107). My core direction is Computational Social Science: I develop agent-based and data-driven models to study digital markets, networked behavior, and policy-relevant outcomes in technology-mediated environments.

Lingnan University’s ambition to strengthen impactful interdisciplinary research in the digital era strongly matches my agenda. My research connects (i) computational modeling of heterogeneous agents and networks, (ii) modern machine learning and AI-assisted research pipelines, and (iii) substantive questions in market dynamics and public policy. In particular, my job-market paper develops an agent-based model of competing ICT firms that adjust pricing and physical service-capacity investment under network effects and congestion, linking firm strategy to welfare and market structure in digital industries.

I began my Ph.D. in 2021, at a time when the COVID-19 pandemic accelerated digitalization and revealed how technological shocks can rapidly reshape markets and economic behavior. 
This experience motivated the central chapter of my dissertation, which studies how firms in ICT markets adjust capacity investment, pricing, and service quality in response to technological innovations, evolving digital infrastructures, and changes in user behavior. 
% Using firm-level data, I estimate key structural parameters and evaluate counterfactual technology-adoption scenarios through simulation-based econometrics, with a focus on how digital technologies influence productivity, competition, and strategic conduct.
My dissertation, \textit{“Strategic Behavior and Economic Dynamics in ICT Markets: An Agent-Based Framework,”} develops computational models in which firms adapt to technological change, heterogeneous users, and congestion effects. 
These models allow me to examine how innovations propagate across markets and how firm strategies evolve under different technological environments. 
To complement this theoretical foundation, I use Monte Carlo simulations and Neural Networks surrogate models to estimate behavioral parameters and assess the economic mechanisms of alternative technological and regulatory scenarios. 

Beyond my dissertation, I have built an interdisciplinary publication pipeline that aligns with Lingnan’s priorities on data science and human-centric technology. First, my peer-reviewed article \emph{Polytope Fraud Theory} (\textit{International Review of Financial Analysis}, 2025; with D.\ Zhao, F.\ Schweizer, and D.\ Sornette) uses supervised learning to identify fraud-related anomalies in public financial data. Second, my coauthored working paper \emph{Agentic Workflows for Economic Research} (arXiv; with H.\ Wang, J.\ Yi, H.\ Dawid, and P.\ Harting) proposes and implements a human-in-the-loop framework using LLM-based agents to improve efficiency and reproducibility across the research lifecycle. Together, these projects demonstrate my capacity to combine conceptual framing with rigorous computational methods and to contribute to a university-wide “digitalized future” in research and teaching.

% \textbf{Research and funding plan.} 
If appointed, I will pursue a coherent 2--3 year research program around digital-economy behavior and policy:
\begin{itemize}
    \item \textbf{Networked markets and platform governance:} how congestion, switching, and network effects interact to shape competition and consumer welfare, with implications for platform regulation and industrial policy;
    \item \textbf{Human--AI interactions in economic measurement:} using agentic/LLM workflows to construct scalable datasets and to benchmark robust, transparent research pipelines;
    \item \textbf{Digital finance and disclosure analytics (selective extension):} anomaly detection and investor-protection questions grounded in transparent ML and reproducible evaluation.
\end{itemize}
I have been trained in a highly international, grant-funded environment (EPOC Marie Sk\l odowska--Curie), and I am prepared to develop competitive proposals in Hong Kong (e.g., early-career and interdisciplinary schemes) with clear deliverables, open-source components, and policy-facing outputs.

% \textbf{Teaching and academic service.} 
I value Lingnan’s whole-person education mission and would contribute enthusiastically to undergraduate and postgraduate teaching. I have experience as a Helping Assistant in quantitative business administration and supported teaching content for \textit{Microeconomics}, \textit{Corporate Finance}, and \textit{Business Administration}. I can teach courses such as:
\begin{itemize}
    \item Computational Social Science / Computational Economics
    \item Data Science for Social Sciences (Python-based)
    \item Industrial Organization and Digital Platform Markets
    \item Machine Learning for Policy and Business Analytics (applied, reproducible workflows)
\end{itemize}
I also contribute actively to academic community-building (e.g., organizing and speaking at workshops on AI agents in science and presenting at international conferences), and I would be keen to collaborate across Lingnan’s Schools and Faculties.

% I learned about this opportunity via \textbf{Computeroxy}. Enclosed please find my CV and supporting materials. My referees (available upon request in the application system) include Prof.\ Herbert Dawid, Prof.\ Didier Sornette, and Prof.\ Domenico Delli Gatti.
Thank you so much for your consideration! I would be delighted to discuss how my work in computational social science and digital-economy research can contribute to Lingnan University’s interdisciplinary priorities and research impact.

\vspace{6mm}
\noindent Sincerely,\\[10mm]
\noindent Zhongli Wang

\end{document}
